\begin{figure}[t]
    \centering
    \begingroup
    % Keep these fills identical to GROUP_BACKGROUNDS in
    % plot_scripts/plot_resolved_tiling_sensitivity.py.
    \definecolor{controlflowbg}{HTML}{F8D7DA}
    \definecolor{addresscalcgbg}{HTML}{D9EAD3}
    \definecolor{mmschedulebg}{HTML}{DDEBF7}
    \newcommand{\controlmath}[1]{{\setlength{\fboxsep}{1.0pt}%
        \colorbox{controlflowbg}{\strut\ensuremath{#1}}}}
    \newcommand{\addressmath}[1]{{\setlength{\fboxsep}{1.0pt}%
        \colorbox{addresscalcgbg}{\strut\ensuremath{#1}}}}
    \newcommand{\schedulemath}[1]{{\setlength{\fboxsep}{1.0pt}%
        \colorbox{mmschedulebg}{\strut\ensuremath{#1}}}}
    \newcommand{\auxfn}[1]{{\sffamily\bfseries #1}}
    \newcommand{\controlgroup}[1]{{\setlength{\fboxsep}{1.5pt}%
        \colorbox{controlflowbg}{\parbox{0.91\linewidth}{#1}}}}
    \newcommand{\addressgroup}[1]{{\setlength{\fboxsep}{1.5pt}%
        \colorbox{addresscalcgbg}{\parbox{0.91\linewidth}{#1}}}}
    \newcommand{\schedulegroup}[1]{{\setlength{\fboxsep}{1.5pt}%
        \colorbox{mmschedulebg}{\parbox{0.91\linewidth}{#1}}}}

    \begin{minipage}[t]{0.405\linewidth}
        \fontsize{5.45}{6.05}\selectfont
        \setlength{\fboxsep}{3pt}
        \fcolorbox{black!30}{white}{%
        \parbox[t][0.72in][t]{0.92\linewidth}{%
            \centering
            \textbf{(a) D-Attn \(\mathrm{Query}\times\mathrm{Key}^{\mathsf T}\)}\par
            \vspace{1pt}
            \(\displaystyle
            \mathrm{Output}_{b,h}
              =\mathrm{Query}_{b,h}\mathrm{Key}_{b,h}^{\mathsf T}
            \)\par
            \vspace{1pt}
            \raggedright
            \(\mathrm{Query}_{b,h}:[1,D_{head}]\): Query row of attention
            head \(h\) for request \(b\) in the batch.\par
            \(\mathrm{Key}_{b,h}:[T_b,D_{head}]\): cached Key rows of
            attention head \(h\) for request \(b\) in the batch.\par
            \(\mathbf{T}=[T_0,\ldots,T_{B-1}]\): per-request KV lengths
            (runtime input).
        }}\par
        \vspace{3pt}
        \fontsize{5.05}{5.50}\selectfont
        \fcolorbox{black!30}{white}{%
        \parbox[t][1.15in][t]{0.92\linewidth}{%
            \centering
            \textbf{(b) Representative schedule parameters for
                \(\mathrm{Query}\times\mathrm{Key}^{\mathsf T}\)}(34 total)\par
            \vspace{2pt}
            \raggedright
            \controlgroup{%
                \centering
                \textbf{Control logic (6 total)}\par
                \raggedright
                \(T_N\): Key rows in one output tile.\\
                \(P\): Key rows in one cache page.\\
                \(T_D\): Key columns in one narrow load.
            }\par
            \vspace{1pt}
            \addressgroup{%
                \centering
                \textbf{Per-core tile dimensions/addressing
                (14 total)}\par
                \raggedright
                \(M_c\): Output-row extent assigned to one core.\\
                \(K_c\): Reduction-column extent assigned to one core.
            }\par
            \vspace{1pt}
            \schedulegroup{%
                \centering
                \textbf{Core-loop scheduling/memory resources
                (14 total)}\par
                \raggedright
                \(N_c\): Number of active compute cores.\\
                \(n_K\): Number of Key buffers.
            }
        }}
    \end{minipage}\hfill
    \begin{minipage}[t]{0.575\linewidth}
        \fontsize{4.80}{5.20}\selectfont
        \setlength{\fboxsep}{3pt}
        \fcolorbox{black!30}{white}{%
        \parbox[t][2.10in][t]{0.94\linewidth}{%
            \centering
            \textbf{(c) Tiled BMM pseudocode for
                \(\mathrm{Query}\times\mathrm{Key}^{\mathsf T}\)}\par
            \vspace{2pt}
            \raggedright
            \ttfamily
            \begin{tabular}{@{}>{\color{black!55}}r@{\hspace{3pt}}l@{}}
             1 & \textbf{function} \auxfn{BMM}(Query, Key, \(\mathbf{T}\);
                 \(\controlmath{T_N},\controlmath{P},\controlmath{T_D},
                 \addressmath{M_c},\addressmath{K_c},
                 \schedulemath{N_c},\schedulemath{n_K}\)) \\
             2 & \quad configure \(\schedulemath{N_c}\) cores and
                 \(\schedulemath{n_K}\) Key buffers \\
             3 & \quad\textbf{parallel for} \(c=0,\ldots,\schedulemath{N_c}-1\) \\
             4 & \quad\quad\textbf{for} tile \(\tau\) with
                 \(\controlmath{T_N}\) Key rows \(\times\)
                 \(\addressmath{K_c}\) columns assigned to core \(c\) \\
             5 & \quad\quad\quad \textit{// \(n_0\): first Key row;
                 \(k_0\): first reduction column; \(T_b\): request \(b\)'s KV length} \\
             6 & \quad\quad\quad
                 \((b,h,n_0,k_0)\leftarrow
                 \text{\auxfn{TileOffset}}(\tau;\addressmath{M_c},
                 \controlmath{T_N},\addressmath{K_c},\mathbf{T})\) \\
             7 & \quad\quad\quad
                 \(n\leftarrow\min(\controlmath{T_N},T_b-n_0)\) \\
             8 & \quad\quad\quad \textbf{for} \((p,r,u)\) in
                 \auxfn{PageFragments}\((\mathrm{Key},b,n_0,n;\controlmath{P})\) \\
             9 & \quad\quad\quad\quad\textbf{if}
                 \(\controlmath{T_N}=\controlmath{P}\) \\
            10 & \quad\quad\quad\quad\quad
                 one-copy load Key\([p,r:r+u,h,k_0:k_0+\addressmath{K_c}]\)
                 \(\rightarrow\mathrm{KeyBuf[next(}\schedulemath{n_K}\mathrm{)]}\) \\
            11 & \quad\quad\quad\quad\textbf{else}
                 \(\;(\controlmath{T_N}\ne\controlmath{P})\) \\
            12 & \quad\quad\quad\quad\quad\textbf{for}
                 \(d_0=k_0,k_0+\controlmath{T_D},\ldots,
                 k_0+\addressmath{K_c}\) \\
            13 & \quad\quad\quad\quad\quad\quad
                 append load Key\([p,r:r+u,h,d_0:d_0+\controlmath{T_D}]\)
                 \(\rightarrow\mathrm{KeyBuf[next(}\schedulemath{n_K}\mathrm{)]}\) \\
            14 & \quad\quad\quad
                 load Query\([b,0,h,k_0:k_0+\addressmath{K_c}]\)
                 \(\rightarrow\mathrm{QueryBuf}\) \\
            15 & \quad\quad\quad Output\(_{\rm tile}\mathrel{+}=\)
                 \auxfn{MMAD}(QueryBuf, KeyBuf)
            \end{tabular}
            \vspace{3pt}\par
            \rmfamily
            \auxfn{TileOffset}: maps tile \(\tau\) to start coordinate
            \((b,h,n_0,k_0)\): request, head, Key row, and reduction column.\par
            \auxfn{PageFragments}: resolves \(n\) logical Key rows through the
            page table as non-crossing fragments \((p,r,u)\).\par
            \auxfn{MMAD}: multiplies QueryBuf by KeyBuf and accumulates the
            partial result into Output\(_{\rm tile}\).
        }}
    \end{minipage}

    \caption{Tiled D-Attn BMM. Panel (b) lists seven representative
    schedule parameters; the three groups contain 6, 14, and 14 schedule
    parameters in total. Matching backgrounds in panel (c) mark their use
    sites and correspond to the groups evaluated in
    \autoref{fig:eval:constant-sensitivity}.}
    \Description{Three panels. Panel (a) defines the per-request Query--Key
    multiplication, explains the Query and Key tensors for request and head
    indices, and defines the runtime array of per-request KV lengths. Panel
    (b) lists seven representative schedule parameters and gives total
    counts of 6, 14, and 14 schedule parameters for Control logic, Per-core tile
    dimensions/addressing, and Core-loop scheduling/memory resources,
    respectively. Panel (c) takes all seven schedule
    parameters as inputs and shows core configuration, per-core tile indexing,
    page-aware Key loading with separate T-N-equals-P and T-N-not-equals-P
    paths, and the final tiled matrix multiplication.}
    \label{fig:resolved-tiling-code}
    \endgroup
\end{figure}
